\section{A nonreal-divisor budget strictly improves the three-factor gate}
\label{sec:s5-nonreal-background}

\paragraph{Review status.}
The complete post-writing author reading and two independent
whole-proof readings have passed. The primary-formula checks
and the unused zero-count source issue have their own recorded scope.

Retain the actual three completed Hecke functions of MA-RB-013,
indexed by \(j\in\{B,T,U\}\), with the exact all-place
conductors, gamma exponents and logarithmic derivatives
\(H_j\) in (MA-RB.200)--(MA-RB.201). Their entire completions,
endpoint nonvanishing, positive integer divisors, reflection,
and positive zero sums are precisely (MA-RB.210)--(MA-RB.211).
In particular, for \(\sigma\ge1\),
\[
 H_j(\sigma)=\sum_{\rho\in\mathcal Z_j}
       \frac{\sigma-\operatorname{Re}\rho}{|\sigma-\rho|^2}>0.
 \tag{MA-RB.250}
\]
All multiplicities are retained. The containing strip is
\(0\le\operatorname{Re}\rho\le1\), and the real endpoints
zero and one are absent. No new zero-location hypothesis is imposed.

Define the following constants and functions from the actual
logarithmic derivatives on the real half-line:
\[
 \begin{aligned}
  C_j&=\min_{u\ge1}(u-\tfrac12)H_j(u),&
  r_j&=\lfloor C_j\rfloor,\\
  h_j(R)&=\frac{H_j(R)}{R-\tfrac12}
                 -\frac{r_j}{R(R-1)}\qquad(R\ge2),\\
  V_j(\sigma)&=\sup_{R\ge\max\{2,\sigma\}}
                              \max\{0,h_j(R)\},\\
  D_j(\sigma)&=(\sigma-\tfrac12)V_j(\sigma),&
  G_j(\sigma)&=H_j(\sigma)-D_j(\sigma).
 \end{aligned}
 \tag{MA-RB.251}
\]
The use of a minimum and the finiteness of these expressions
are proved below. The integer \(r_j\) is an upper bound,
not a claim to know the actual number of real zeros.

\paragraph{MA-RB-014 (one all-parameter nonreal-background refinement).}
For each \(j\), let \(n_j^{\mathrm{re}}\) be the total multiplicity
of the real zeros of \(\Lambda_j\). Then
\[
       n_j^{\mathrm{re}}\le r_j<\infty.
 \tag{MA-RB.252}
\]
The supremum in (MA-RB.251) is a positive attained maximum
for every \(\sigma\ge1\). Both \(D_j\) and \(G_j\) are
continuous there, and \(D_j\) is a lower bound for the
nonreal part of (MA-RB.250). If \(r_j\ge1\), then
\[
                    0<D_j(\sigma)<H_j(\sigma),
                    \qquad G_j(\sigma)>0.
 \tag{MA-RB.253}
\]

For an integer \(k\ge1\), if any \(r_j<2k\), the residual
comparison branch of order \(k\) in MA-RB-002 is impossible.
Otherwise define
\[
 \begin{aligned}
  \widehat F_{j,k}(\sigma)
       &=(\sigma-\tfrac12)^2
                         -\frac{k(2\sigma-1)}{G_j(\sigma)},\\
  \widehat{\mathfrak K}_{L,k}
       &=\min_{j\in\{B,T,U\}}\inf_{\sigma>1}
                                      \widehat F_{j,k}(\sigma).
 \end{aligned}
 \tag{MA-RB.254}
\]
These infima are finite and attained after adjoining
\(\sigma=1\). For the excluded case, put
\(\widehat{\mathfrak K}_{L,k}=-\infty\).
With this convention the strict comparison is unconditional:
\[
 \widehat{\mathfrak K}_{L,k}
          <\mathfrak K_{L,k}<\mathfrak J_{L,k}<\frac14.
 \tag{MA-RB.255}
\]
Here the middle two bounds are the complete envelopes of
MA-RB-013 and MA-RB-012.

If the actual comparison fails with real abscissa
\(\gamma\in(\tfrac12,1)\) and integer pole order \(k\), then
all \(r_j\ge2k\) and
\[
 \begin{gathered}
   (\gamma-\tfrac12)^2\le\widehat{\mathfrak K}_{L,k}
                         \le\widehat{\mathfrak K}_{L,1},\\
   \operatorname{ord}_{\gamma}\Lambda_j
     =\operatorname{ord}_{1-\gamma}\Lambda_j\ge k
                                  \quad(j=B,T,U).
 \end{gathered}
 \tag{MA-RB.256}
\]
Thus \(\widehat{\mathfrak K}_{L,1}\le0\), including the
excluded case, suffices for equality of the comparison.
If it is positive, the existing sufficient shared-zero gate
is restricted further to
\[
 B(\gamma)=L_U(\gamma)=0,\qquad
 \frac12<\gamma\le\frac12+
                    \sqrt{\widehat{\mathfrak K}_{L,1}}<1,
 \tag{MA-RB.257}
\]
with \(L_T(\gamma)=0\) an additional necessary condition.
The theorem does not assert that this final interval is empty.

\begin{proof}
\emph{The real multiplicity bound and its minimum.}
The real zero set is finite, since it lies in the compact
interval \([0,1]\), the completion is entire and nonzero,
and its two endpoints are nonzero. A real reflection pair
\(\beta,1-\beta\), with
\(v=\beta-\tfrac12\) and \(x=u-\tfrac12\), contributes
\[
 \frac1{u-\beta}+\frac1{u-1+\beta}
        =\frac{2x}{x^2-v^2}\ge\frac2x.
 \tag{MA-RB.258}
\]
A zero at \(\beta=\tfrac12\) is counted with its actual
multiplicity, equivalently with half the paired weight.
Positivity of every other zero term gives
\(n_j^{\mathrm{re}}\le(u-\tfrac12)H_j(u)\) for every
\(u\ge1\).

For later use record the precise real-axis asymptotics.
Writing \(d_j=d_j^++d_j^-\), the classical digamma and
trigamma expansions, DLMF 5.11.2 and 5.15.8, give from
(MA-RB.201)
\[
 \begin{aligned}
 H_j(R)&=\frac{d_j}{2}\log R+c_j+O_L(R^{-1}),\\
 H_j'(R)&=\frac{d_j}{2R}+O_L(R^{-2}),\qquad
 c_j=\frac12\log N_j-\frac{d_j}{2}\log(2\pi).
 \end{aligned}
 \tag{MA-RB.259}
\]
The Euler contribution and its first derivative are
\(O_L(2^{-R})\): the actual inertia-invariant coefficient
bound in (MA-RB.240) bounds them respectively by constants
times \(\sum_{n\ge2}(\log n)n^{-R}\) and
\(\sum_{n\ge2}(\log n)^2n^{-R}\). For \(R\ge2\),
factor \(n^{-R}=n^{-2}n^{-(R-2)}\) to obtain the asserted
exponential bounds. Thus no differentiation of an unsupported
big-O remainder is being used.

The continuous positive function
\((u-\tfrac12)H_j(u)\) tends to infinity by (MA-RB.259).
It therefore attains a finite positive minimum on
\([1,\infty)\). Taking the minimum in (MA-RB.258), followed
by the integer part, proves (MA-RB.252).

\emph{The exact comparison of reflection pairs.}
Fix \(\sigma\ge1\) and \(R\ge\max\{2,\sigma\}\).
Put
\[
 x=\sigma-\tfrac12,\quad X=R-\tfrac12,\quad
 u=\operatorname{Re}\rho-\tfrac12,\quad
 y=(\operatorname{Im}\rho)^2,
 \qquad A=x^2+y-u^2,\quad B=X^2+y-u^2.
 \tag{MA-RB.260}
\]
The contribution of a pair under
\(\rho\mapsto1-\overline\rho\), divided by \(2x\),
is \(A/(A^2+4u^2y)\), as in (MA-RB.214).
At the point \(R\), the same normalized contribution
is \(B/(B^2+4u^2y)\). Their exact difference is
\[
 \frac{A}{A^2+4u^2y}-\frac{B}{B^2+4u^2y}
   =\frac{(B-A)(AB-4u^2y)}
          {(A^2+4u^2y)(B^2+4u^2y)}.
 \tag{MA-RB.261}
\]
Here
\[
 AB-4u^2y
   =y^2+(x^2+X^2-6u^2)y
                     +(x^2-u^2)(X^2-u^2)\ge0,
 \tag{MA-RB.262}
\]
because \(x^2\ge1/4\), \(X^2\ge9/4\), and
\(u^2\le1/4\). In fact the coefficient of \(y\)
is at least one. For a nonreal pair, \(y>0\), so this
last expression is strictly positive. All denominators
are positive; in the real case the excluded endpoints
remove the only possible singularity at \(\sigma=1\).
Consequently the normalized comparison is nonnegative,
and it is strict for every nonreal pair if \(R>\sigma\).
Fixed points of the reflection are again treated with
half the paired weight. Absolute convergence and positivity
allow summation without any ordering qualification.

Write \(H_j^{\mathrm{re}}\) and \(H_j^{\mathrm{nr}}\)
for the two parts of (MA-RB.250). The result just proved is
\[
 \frac{H_j^{\mathrm{nr}}(\sigma)}{\sigma-\tfrac12}
       \ge\frac{H_j^{\mathrm{nr}}(R)}{R-\tfrac12}.
 \tag{MA-RB.263}
\]
If the nonreal divisor is nonempty and \(R>\sigma\),
the inequality is strict. On the other hand, a real pair
at \(R>1\) is at most
\(2(R-\tfrac12)/[R(R-1)]\), by \(u^2\le1/4\).
Thus (MA-RB.252) gives
\[
 0\le\frac{H_j^{\mathrm{re}}(R)}{R-\tfrac12}
       \le\frac{n_j^{\mathrm{re}}}{R(R-1)}
       \le\frac{r_j}{R(R-1)}.
 \tag{MA-RB.264}
\]
Combining (MA-RB.263)--(MA-RB.264) yields
\[
 \frac{H_j^{\mathrm{nr}}(\sigma)}{\sigma-\tfrac12}
                    \ge\max\{0,h_j(R)\}.
 \tag{MA-RB.265}
\]

\emph{Attainment, continuity and the strict denominator.}
By (MA-RB.259),
\[
 h_j(R)=\frac{d_j\log R}{2R}+O_L(R^{-1})>0
       \quad\hbox{for all sufficiently large }R,
       \qquad h_j(R)\longrightarrow0.
 \tag{MA-RB.266}
\]
The function \(h_j\) is continuous on \([2,\infty)\).
For any fixed lower endpoint \(a\ge2\), choose
\(R_1\ge a\) with \(h_j(R_1)>0\). Beyond a sufficiently
large finite endpoint the values of \(h_j\) have smaller
absolute value than \(h_j(R_1)/2\). Its positive supremum
on \([a,\infty)\) is therefore attained in a compact
interval. In particular the positive part in (MA-RB.251)
does not conceal a nonattained value zero.

The same argument proves continuity of the running maximum.
Near a fixed \(a_0\ge2\), choose \(R_1>a_0\) with
\(h_j(R_1)>0\), and choose one finite upper endpoint
that works for all nearby \(a\). The maximum is then
\(\max_{R\in[a,T]}h_j(R)\). Uniform continuity of \(h_j\)
on the containing compact interval shows that this maximum
depends continuously on \(a\). Composing with
\(a=\max\{2,\sigma\}\) proves the claimed continuity
of \(V_j,D_j,G_j\).

Equation (MA-RB.265) now proves
\(0<D_j(\sigma)\le H_j^{\mathrm{nr}}(\sigma)\).
It also proves that the nonreal divisor is nonempty,
without an external zero-count estimate. Let \(R_*\)
be an actual maximizing point in (MA-RB.251).
If \(R_*>\sigma\), the strict form of (MA-RB.263) gives
\(D_j(\sigma)<H_j^{\mathrm{nr}}(\sigma)\le H_j(\sigma)\).
If \(R_*=\sigma\), necessarily \(\sigma\ge2\), and
the defining identity gives instead
\[
 G_j(\sigma)=\frac{r_j(\sigma-\tfrac12)}{\sigma(\sigma-1)}.
 \tag{MA-RB.267}
\]
This is strictly positive whenever \(r_j\ge1\).
These two cases prove (MA-RB.253), including the endpoint
\(\sigma=1\), where every maximizing point is strictly larger.

\emph{The tail and existence of the new minima.}
Using both asymptotics in (MA-RB.259), direct differentiation
of the definition of \(h_j\) gives
\[
 \begin{aligned}
 h_j'(R)
  &=\frac{(R-\tfrac12)H_j'(R)-H_j(R)}{(R-\tfrac12)^2}
       +\frac{r_j(2R-1)}{R^2(R-1)^2}\\
  &=-\frac{d_j\log R}{2R^2}+O_{L,r_j}(R^{-2})<0
                    \quad\hbox{for all sufficiently large }R.
 \end{aligned}
 \tag{MA-RB.268}
\]
Hence, for every sufficiently large \(\sigma\), the running
maximum is achieved at \(R=\sigma\). Equation (MA-RB.267)
then holds identically on that tail. If \(r_j\ge2k\),
\[
 \widehat F_{j,k}(\sigma)
   =\frac14+\sigma(\sigma-1)
                           \left(1-\frac{2k}{r_j}\right)
   \ge\frac14
               \quad\hbox{for all sufficiently large }\sigma.
 \tag{MA-RB.269}
\]
When \(r_j>2k\) this tends to infinity; when \(r_j=2k\)
it is exactly \(1/4\) there. On the other hand,
\(\widehat F_{j,k}(1)=1/4-k/G_j(1)<1/4\).
Continuity therefore puts each minimum in a compact interval
of \([1,\infty)\). It is finite and attained there.
This argument avoids an incorrect coercivity assertion
in the equality case \(r_j=2k\).

For every \(\sigma\ge1\), (MA-RB.253) implies
\[
   \widehat F_{j,k}(\sigma)<F_{j,k}(\sigma),
 \tag{MA-RB.270}
\]
where \(F_{j,k}\) is the old function in (MA-RB.202).
Choose an actual minimizing pair \((j_*,\sigma_*)\)
for the old three-factor envelope, allowing \(\sigma_*=1\).
Evaluate (MA-RB.270) there; the continuous endpoint value
is admitted by the infimum over \(\sigma>1\).
This proves the first strict inequality in (MA-RB.255)
whenever all \(r_j\ge2k\). The excluded case satisfies
it by the stipulated value \(-\infty\).

\emph{Transfer to the actual residual comparison.}
If comparison fails with the actual pair \((\gamma,k)\),
(MA-RB.238) gives zeros of multiplicity at least \(k\)
at each of the distinct real points \(\gamma,1-\gamma\)
in every \(\Lambda_j\). Therefore
\(2k\le n_j^{\mathrm{re}}\le r_j\), proving the branch
test. These two real zeros are disjoint from the nonreal
divisor used in (MA-RB.265). Keeping both contributions
in the same positive sum gives
\[
 H_j(\sigma)\ge D_j(\sigma)
       +\frac{k(2\sigma-1)}
          {(\sigma-\tfrac12)^2-(\gamma-\tfrac12)^2}.
 \tag{MA-RB.271}
\]
The denominator is positive for \(\sigma\ge1\).
Subtract \(D_j\), use \(G_j>0\), and rearrange to obtain
\((\gamma-\tfrac12)^2\le\widehat F_{j,k}(\sigma)\).
Taking the three infima proves the first part of (MA-RB.256).
The background budget is independent of \(k\); hence these
functions and their infima decrease with \(k\) throughout
the feasible range. Once a branch is excluded, the value
\(-\infty\) preserves this monotonicity. This proves the
remaining assertions.
\end{proof}

\paragraph{Finite upper certificates without a zero search.}
All the quantities used in the strict comparison can be
approximated in the direction required for a rigorous upper
certificate. At any fixed rational \(u>1\), use the full
finite Euler sum and the explicit tail (MA-RB.240) to produce
upper and lower values \(H_j^+(u),H_j^-(u)\) converging
to \(H_j(u)\). The known conductor and gamma terms are
included in both. Taking finite minima of
\((u-\tfrac12)H_j^+(u)\), over rational points and increasing
truncations, gives upper approximants \(C_j^+\downarrow C_j\).
Continuity at one allows the minimizing point to be one;
no Euler series at one is required. Thus
\[
 r_j^+=\lfloor C_j^+\rfloor\ge r_j,
 \qquad r_j^+=r_j\quad\hbox{eventually}.
 \tag{MA-RB.272}
\]
The eventual equality also holds if \(C_j\) is an integer:
eventually \(C_j\le C_j^+<C_j+1\). This is a convergence
statement, not an effective stopping test for equality.
Any finite certificate \(r_j^+<2k\) already excludes the
actual order-\(k\) branch.

Otherwise fix a rational \(\sigma>1\) and finitely many
rational \(R_\ell\ge\max\{2,\sigma\}\), and put
\[
 \begin{aligned}
  V_j^-(\sigma)
   &=\max\left(0,\max_\ell
       \left\{\frac{H_j^-(R_\ell)}{R_\ell-\tfrac12}
                       -\frac{r_j^+}{R_\ell(R_\ell-1)}\right\}\right),\\
  G_j^+(\sigma)
   &=H_j^+(\sigma)-(\sigma-\tfrac12)V_j^-(\sigma).
 \end{aligned}
 \tag{MA-RB.273}
\]
When the exact branch is feasible, these satisfy
\(0<G_j(\sigma)\le G_j^+(\sigma)\), and therefore
\[
 \widehat F_{j,k}(\sigma)
    \le(\sigma-\tfrac12)^2
                     -\frac{k(2\sigma-1)}{G_j^+(\sigma)}.
 \tag{MA-RB.274}
\]
If instead the exact branch is excluded, the right side of
(MA-RB.274), whenever its computed denominator is positive,
still bounds the extended global value \(-\infty\); no
undefined individual function is used in that case.
As the finite lists of rational \(u,R_\ell,\sigma\)
become dense and the Euler truncations increase, the upper
certificates recover the feasible infimum. Indeed, (MA-RB.272)
eventually gives the exact integer, the attained maximum in
(MA-RB.251) is approximated by rational points, and continuity
then gives convergence at each fixed \(\sigma\). Taking finite
minima over rational \(\sigma>1\) also reaches an endpoint
minimum at one by a limit. No effective stabilization of
\(r_j\), no uniform Euler convergence at one, and no actual
zero enumeration is asserted.

\paragraph{Source, finite-cover and arithmetic scope.}
The Hecke functions and Artin objects are exactly the induced
rank-one local systems of (MA-RB.206)--(MA-RB.210), with the
finite-cover conventions already recorded from SGA~1,
Expos\'e~V, \S7, and Deligne (1979), Definition~6.1 and
\S\S6.2--6.3. This argument only uses their established
completed divisors and all-place Euler bounds. DLMF 5.11.2
and 5.15.8 are cited for the classical real-axis asymptotics;
the paired inequalities, running maximum, integer branch
test, denominator argument and strict infimum comparison
are proved here. No explicit zero-count theorem or numerical
zero data enters this proof. In particular, no root-number
parity refinement is assumed.

The result makes the pre-existing real shared-zero interval
strictly smaller, or excludes it if the multiplicity test
already succeeds. It does not establish full comparison,
\((\mathcal R_B)\), or the absence of a remaining shared
zero in a case where the new upper bound is positive.
